\graphicspath{{lectures/08_eloadas_abrak/}}

\section{8. előadás -- Állandó együtthatós n-edrendű lineáris differenciálegyenletek}

\subsection{Állandó együtthatós n-edrendű differenciálegyenlet}
\textit{EA: 2019. 11. 13. (8).}

Az egyenlet:
\[
 y^{(n)}+a_{n-1}y^{(n-1)}+\cdots+a_1y'+a_0y=b,
\]
ahol
\[
 a_i\in\mathbb R\quad(i=0,\ldots,n-1),\qquad
 b:I\longrightarrow\mathbb R\quad\text{folytonos}.
\]

\textbf{A homogén egyenletnek} van $n$ darab lineárisan független
megoldása.

\textbf{Előállítás.} Ötlet: keressünk megoldást
\[
 \varphi(t)=e^{\lambda t}
\]
alakban. Ekkor
\[
 \varphi^{(k)}(t)=\lambda^k e^{\lambda t},
 \qquad k=1,\ldots,n,
\]
és a homogén egyenletbe helyettesítve:
\[
 \bigl(\underbrace{\lambda^n+a_{n-1}\lambda^{n-1}
              +\cdots+a_1\lambda+a_0}_{K(\lambda)}\bigr)
 e^{\lambda t}=0.
\]
A $K(\lambda)$ a \textbf{karakterisztikus polinom}
(a homogén egyenleté).

\subsection{Komplex értékű differenciálegyenletek}
\textit{(Állandó együtthatós, lineáris egyenletek.)}

Tekintsük a komplex értékű függvényt:
\[
 \varphi:\mathbb R\longrightarrow\mathbb C,\qquad
 \varphi(t)=\varphi_1(t)+i\,\varphi_2(t)
 \quad(t\in\mathbb R),
\]
ahol
\[
 \varphi_1,\varphi_2:\mathbb R\longrightarrow\mathbb R,\qquad
 \varphi_1(t)=\operatorname{Re}(\varphi(t)),\qquad
 \varphi_2(t)=\operatorname{Im}(\varphi(t)).
\]
A differenciálhatóságra a kézirat jelölésével:
\[
 \varphi\in D\quad\Longleftrightarrow\quad
 \varphi_1,\varphi_2\in D,
 \qquad \varphi'=\varphi_1'+i\,\varphi_2'.
\]
Tehát komplex értékű megoldásokat is lehet értelmezni.

\subsubsection*{Komplex n-edrendű homogén lineáris differenciálegyenlet}
\[
 \underbrace{L_n(y)}_{\text{differenciáloperátor}}
 :=y^{(n)}+a_{n-1}y^{(n-1)}+\cdots+a_1y'+a_0y=0,
 \qquad a_i\in\mathbb C\quad(i=0,\ldots,n-1).
 \tag{KH}
\]
Komplex megoldások értelmezhetők.
Az inhomogén komplex megoldásaira a valós esethez hasonló eredmények
érvényesek.

\textbf{Új probléma:} valós együttható esetben a valós megoldások
kiválasztása.

\textbf{1. megjegyzés.} Ha $\lambda\in\mathbb C$ és
$\lambda=\alpha+i\beta$ ($\alpha,\beta\in\mathbb R$), akkor az
Euler-formula:
\[
 \boxed{e^{it}=\cos t+i\sin t\qquad(t\in\mathbb R).}
\]
\textbf{2. megjegyzés.} Komplex függvénytan: Cauchy-féle alaptétel stb.
\textit{(Egyetemi jegyzet: Szőkefalvi-Nagy Béla: Komplex függvénytan.)}

\textbf{3. megjegyzés.}
\begin{enumerate}[label=\roman*)]
\item Ha $\lambda\in\mathbb C$, $\lambda=\alpha+i\beta$
($\alpha,\beta\in\mathbb R$), akkor
\[
 e^{\lambda t}=e^{(\alpha+i\beta)t}
 =e^{\alpha t}\bigl(\cos(\beta t)+i\sin(\beta t)\bigr).
\]
\item Igazolható:
\[
 (e^{\lambda t})'=\lambda e^{\lambda t}.
\]
\end{enumerate}

\begin{matbox}
\textbf{Tétel.} A $\varphi:\mathbb R\longrightarrow\mathbb C$,
$\varphi(t)=e^{\lambda t}$ ($t\in\mathbb R$, $\lambda\in\mathbb C$)
függvény megoldása a komplex homogén egyenletnek pontosan akkor,
ha $\lambda$ gyöke a karakterisztikus polinomnak, azaz
\[
 K(\lambda)=\lambda^n+a_{n-1}\lambda^{n-1}
 +\cdots+a_1\lambda+a_0=0.
\]
\end{matbox}

\subsection{A homogén egyenlet alaprendszere}
\begin{matbox}
\textbf{Tétel.} Tegyük fel, hogy $K(\lambda)$ gyökei különbözők:
\[
 \lambda_1,\lambda_2,\ldots,\lambda_n\in\mathbb C.
\]
Ekkor
\[
 \varphi_k(t)=e^{\lambda_k t}
 \qquad(t\in\mathbb R,\ k=1,2,\ldots,n)
\]
a komplex homogén egyenlet egy alaprendszere.
\end{matbox}

\textbf{Bizonyítás (vázlat).}
\begin{itemize}
 \item Megoldások: igen.
 \item Lineárisan függetlenek (megengedhető).
\end{itemize}
\hfill$\square$

\subsubsection*{Kérdés: többszörös gyökök esete}
\textbf{Ötlet.} Tegyük fel, hogy $\lambda$ kétszeres gyök.
Egy másik megoldás keresése:
\[
 \text{ha }\lambda_1\ne\lambda_2,\qquad
 \frac{e^{\lambda_1 t}-e^{\lambda_2 t}}{\lambda_1-\lambda_2}
 \xrightarrow[\lambda_2\to\lambda_1]{} t\,e^{\lambda_1t}.
\]
Igazolható, hogy ha $\lambda$ kétszeres gyök,
$e^{\lambda t}$ és $t\,e^{\lambda t}$ megoldások.
Sőt, $t^k e^{\lambda t}$ is megoldás.

\begin{matbox}
\textbf{Tétel.} Legyen $\lambda_0\in\mathbb C$ a
$K(\lambda)$ karakterisztikus polinom pontosan $m$-szeres gyöke.
Ekkor
\[
 \begin{aligned}
 \varphi_1(t)&=e^{\lambda_0t},\\
 \varphi_2(t)&=t\,e^{\lambda_0t},\\
 &\ \vdots\\
 \varphi_m(t)&=t^{m-1}e^{\lambda_0t}
 \end{aligned}
 \qquad(t\in\mathbb R)
\]
a komplex homogén differenciálegyenlet $m$ darab megoldása.
% A kézirat 29. oldalán a „lin. fgt.” szöveg át van húzva;
% a függetlenségre a 30. oldal külön visszatér.
\end{matbox}

\textbf{Bizonyítás (vázlat).}

Legyen
\[
 t^j e^{\lambda_0t},\qquad j=0,1,\ldots,m-1.
\]
\textbf{$j=0$:}
\[
 L_n(e^{\lambda_0t})=K(\lambda_0)e^{\lambda_0t}=0.
\]
$\lambda_0$ a $K(\lambda)$ polinomnak $m$-szeres gyöke pontosan akkor,
ha a kézirat szerint
\[
 K(\lambda_0)=K'(\lambda_0)=\cdots=K^{(m-1)}(\lambda_0).
\]
% FORRÁSHŰ: az eredeti 30. PDF-oldalon a feltételsor végén nem szerepel =0.

Legyen $j=1,2,\ldots,m-1$ rögzített.
\textbf{Ötlet:} felcserélhető a paraméter szerinti deriválás
és az $L_n$ differenciáloperátor:
\begin{align*}
 L_n\bigl(t^j e^{\lambda_0t}\bigr)
 &=L_n\!\left(\frac{\partial^j}{\partial\lambda_0}
                e^{\lambda_0t}\right)\\
 &=\frac{\partial^j}{\partial\lambda_0}
         L_n\bigl(e^{\lambda_0t}\bigr)\\
 &=\frac{\partial^j}{\partial\lambda_0}
         \bigl(K(\lambda_0)e^{\lambda_0t}\bigr)\\
 &=K^{(j)}(\lambda_0)e^{\lambda_0t}
   +\binom{j}{1}K^{(j-1)}(\lambda_0)t\,e^{\lambda_0t}
   +\cdots\\
 &\quad+\binom{j}{j}K(\lambda_0)t^j e^{\lambda_0t}
   =0.
\end{align*}
% FORRÁSHŰ: a 30. oldalon a parciális differenciálás nevezője
% a több helyen is írt \partial\lambda_0 (j-edik hatvány nélkül).
\textit{A negyedik egyenlőségnél: Leibniz-szabály.}
\hfill$\square$

A megoldások lineáris függetlensége: \textbf{indirekt}.

\textbf{1. megjegyzés.} Előállítható $n$ darab lineárisan független
megoldás (ha a $K(\lambda)$ gyökei ismertek).

\subsection{Valós megoldások kiválasztása}
\textbf{2. megjegyzés.} Valós megoldások kiválasztása.

\begin{matbox}
\textbf{Tétel.} Ha a komplex homogén lineáris differenciálegyenletben
\[
 \forall i=1,\ldots,n-1:\quad a_i\in\mathbb R,
\]
és $\lambda\in\mathbb C\setminus\mathbb R$ az $K(\lambda)$
polinom $m$-szeres gyöke, akkor $\overline\lambda$ is
$m$-szeres gyöke $K(\lambda)$-nak.
\end{matbox}
% FORRÁSHŰ: a kéziratban a valós együtthatók feltétele i=1,...,n-1,
% és emiatt a_0 nincs explicit megnevezve. Lásd ellenorzes.md.

\textbf{A lehetséges esetek:}
\begin{enumerate}[label=\arabic*$^\circ$)]
\item $\lambda\in\mathbb R$ többszörös gyök:
      valós megoldások.
\item $\lambda=\alpha+i\beta$ ($\beta\ne0$) gyöke a
      $K(\lambda)$-nak, ekkor $\overline\lambda=\alpha-\beta i$
      is gyöke a $K(\lambda)$-nak.
      \[
       \varphi_1(t)=e^{\lambda t}
       =e^{\alpha t}\bigl(\cos(\beta t)+i\sin(\beta t)\bigr),
      \]
      \[
       \varphi_2(t)=e^{\overline\lambda t}
       =e^{\alpha t}\bigl(\cos(\beta t)-i\sin(\beta t)\bigr).
      \]
      Tehát
      \[
       \boxed{e^{\alpha t}\cos(\beta t),\qquad
              e^{\alpha t}\sin(\beta t)}
      \]
      lineárisan független valós megoldások.
\item Ha $\lambda=\alpha+i\beta$ többszörös gyök:
      hasonlóan kiválaszthatók a megoldások.
\end{enumerate}
\textit{Kéziratos hivatkozás: 08\_An\_alk\_2022\_osz.pdf,
16--18. oldal.}

\subsection{n-edrendű állandó együtthatós inhomogén differenciálegyenlet}
A megoldáshoz:
\begin{itemize}
 \item A homogén egyenlet alaprendszere.
 \item Az inhomogén egyenlet partikuláris megoldása:
       \textbf{állandók variálása} (nehéz lehet).
\end{itemize}
\textit{Példa-hivatkozás a kéziratból: Honlapomon $\to$
Segédanyagok $\to$ Differenciálegyenletek, 43--46. oldal.}

\subsubsection*{Az állandók variálása helyett: szuperpozíciós elv}
\begin{matbox}
\textbf{Tétel (szuperpozíciós elv).}
Tegyük fel, hogy $\psi_1,\psi_2$ megoldásai a
\[
 y^n+a_{n-1}y^{(n-1)}+\cdots+a_1y'+a_0y=b_i,
 \qquad b_i:I\longrightarrow\mathbb R\ \text{folytonos},
 \quad i=1,2
\]
egyenleteknek. Ekkor $\psi_1+\psi_2$ megoldása a
\[
 y^n+a_{n-1}y^{(n-1)}+\cdots+a_1y'+a_0y=b_1+b_2
\]
egyenletnek. $\square$
\end{matbox}
% FORRÁSHŰ: a kézirat 31. oldalán itt a vezető tag y^n,
% miközben a 27. oldalon egyértelműen y^{(n)} szerepel.

\subsubsection*{Próbafüggvény-módszer (bevezetés)}
Ha a „jobb oldal”
\[
 P(t)\,e^{\alpha t}
 \bigl(\cos(\beta t)+i\sin(\beta t)\bigr)
\]
alakú, akkor a „bal oldal” esetén is várható ehhez hasonló megoldás.

% ELŐADÁSHATÁR: az EA: 2019. 11. 20. (9) cím utáni levezetés
% már a 9. előadás része, ebben a modulban nem szerepel.
